\section{Introduction}
\label{sec:intro}

Retrieval-augmented generation (RAG)~\citep{lewis2020rag} now powers
deployed systems across medical QA, legal research, and customer
support, yet many papers still evaluate faithfulness with a narrow
template: one retriever, one generator, one scorer, one threshold, and
one benchmark. That template is convenient, but it
\emph{under-identifies} the claim. A result that looks decisive in one
generator/retriever/scorer cell can change sign, shrink, or become a
deployment tradeoff under an equally defensible audit choice.

The trigger for this paper is a refinement-paradox result from our own
prior work: improving per-passage retrieval appeared to reduce answer
faithfulness. The headline was interesting enough to audit, but it does
not survive the controls reported here. At scale, the original
SQuAD/Mistral contrast shrinks to $0.661 / 0.650 / 0.661$ across
baseline / HCPC-v$1$ / HCPC-v$2$; a matched-similarity CCS intervention
is null; threshold-transfer cells change sign; and a cost-aware
head-to-head produces no dominating system.

The audit reframes the original headline as a stronger
methodological result. The problem is not that a coherence hypothesis
``failed''; the problem is that single-cell RAG faithfulness numbers
hide too many degrees of freedom. Under fixed Mistral generations,
three plausible scorers produce raw score contrasts $0.011$, $0.032$,
and $0.140$; z-scored contrasts on the same rows remain separated
($0.071$, $0.231$, $0.337$ SD); a compact Qwen$2.5$ replication
preserves the raw-contrast ordering; and a
two-rater human calibration shows that no automated scorer should be
treated as an unqualified reference standard. The empirical results are not
independent anecdotes: each one corresponds to an axis that a RAG
faithfulness claim must disclose.

The positive contribution is \textbf{ControlledRAG}: a reusable audit
protocol and prescriptive minimum reporting standard for RAG
faithfulness claims. ControlledRAG is not a new retriever, not a new
faithfulness metric, and not a method leaderboard. It asks authors to
vary or explicitly hold fixed seven axes---generator, retriever,
context structure, scorer, human calibration, threshold transfer, and
cost---so reviewers can tell whether a claim is robust or merely a
metric, retriever, generator, threshold, or cost artifact.

\paragraph{Recommended minimum reporting standard for RAG faithfulness
claims.}
\begin{enumerate}
  \item \textbf{Generator.} Family, parameter count, decoding settings,
        and whether outputs are fixed before scoring.
  \item \textbf{Retriever.} Embedding model, indexer, reranker (if
        any), and any post-retrieval refinement step.
  \item \textbf{Context-set structure.} At least two diagnostics beyond
        per-passage similarity (e.g., CCS, answer-span presence, mean
        query similarity).
  \item \textbf{Scorer.} At least two faithfulness scorers, with a
        reported Pearson or Spearman cross-scorer correlation.
  \item \textbf{Human calibration.} A multi-rater panel with inter-rater
        agreement and Spearman correlation against each automated scorer.
  \item \textbf{Threshold transfer.} Sensitivity of any tuned threshold
        $\tau$ across datasets, reported as a transfer table.
  \item \textbf{Cost.} Per-condition latency, indexing time, and
        external-call cost---including baselines.
\end{enumerate}

\paragraph{Contributions.} We make five contributions, each mapped to
the standard.
\begin{enumerate}
  \item \textbf{ControlledRAG standard
    (\S\ref{sec:controlledrag}).} A seven-axis minimum reporting
    standard for RAG faithfulness claims.
  \item \textbf{Fixed-generation metric fragility
    (\S\ref{sec:metric_fragility}).} A $7{,}500$-row Mistral audit, a
    compact Qwen$2.5$ replication, and a $99$-item human calibration
    showing why scorer choice is an audit axis rather than an
    implementation detail.
  \item \textbf{Matched-context diagnostics
    (\S\ref{sec:matched_ccs}).} A pre-registered CCS null plus an
    answer-span/control diagnostic showing that answer-bearing evidence
    and context structure carry complementary signal in the matched
    cell.
  \item \textbf{Transfer and deployment audit
    (\S\ref{sec:threshold_transfer}--\S\ref{sec:cost}).}
    Threshold-transfer, noise, cost-aware baseline, stronger-retriever,
    and long-form stress results that expose the remaining audit axes
    without turning the paper into a method leaderboard.
  \item \textbf{Review artifact.} An anonymized package containing
    auditable per-query outputs, matched context pairs, source tables,
    scripts, supplement traces, and documentation is provided for
    review; public release links will follow acceptance.
\end{enumerate}

\paragraph{Roadmap.} \S\ref{sec:bg} positions the work against prior
RAG evaluation. \S\ref{sec:setup} describes the pipeline and the
audited cells. \S\ref{sec:controlledrag} defines ControlledRAG.
\S\ref{sec:matched_ccs}--\S\ref{sec:cost} run the seven-axis audit.
\S\ref{sec:discussion} discusses what survives and what does not.
\S\ref{sec:related} returns to the broader literature.
\S\ref{sec:limits} states limitations and ethical considerations.
\S\ref{sec:conclusion} concludes.

\section{Background and Related Work}
\label{sec:bg}

\paragraph{The standard RAG pipeline.} A modern RAG
system~\citep{lewis2020rag,gao2024ragsurvey} threads a dense retriever
\citep{karpukhin2020dpr,xiong2024embedder} returning the top-$k$
passages by query--passage similarity, an optional cross-encoder
reranker~\citep{nogueira2019passage}, an optional refinement module
that further filters or compresses each retained passage
(e.g., extractive sentence selection, query-conditioned summarisation,
or learned compression~\citep{xu2024recomp}), and a generator that
produces the final answer conditioned on the assembled context. Each
stage individually is reported to improve task performance on standard
benchmarks~\citep{karpukhin2020dpr,izacard2022few,xiong2024embedder};
the implicit assumption is that better local retrieval produces more
faithful answers globally.

\paragraph{Decoder-side correction.} Two influential lines move the
correction \emph{into} the decoder.
Self-RAG~\citep{asai2024selfrag} fine-tunes the generator to emit
``reflection tokens'' that decide whether to retrieve, retrieve again,
or stop, paying a fine-tune-time cost and a $2$--$3\times$ inference
latency cost. CRAG~\citep{yan2024crag} introduces a lightweight
retrieval evaluator that triggers web search when document relevance is
judged poor, paying an external-call latency cost. Both take the
position that the generator (or a learned auxiliary) is the right place
to fix retrieval failures. ControlledRAG is orthogonal: it does not
prescribe a place to fix retrieval and does not declare a winner among
deployment policies. It prescribes how a faithfulness claim should be
\emph{reported} so that a reader can tell from the paper which
deployment policy would actually have looked better under a different
scorer or retriever.

\paragraph{Hierarchical retrieval.} RAPTOR~\citep{raptor2024} clusters
retrieved passages into a tree and retrieves from both leaves and
LLM-summarised internal nodes, producing higher-coherence summary
contexts at the cost of an offline LLM-summarisation pass. We compare
against a $2$-level RAPTOR baseline (\S\ref{sec:cost}) and find no
single dominating system across faithfulness, hallucination, and
latency.

\paragraph{Faithfulness evaluation.} Faithfulness---whether a generated
answer is supported by the retrieved context---is typically assessed
via NLI-style scoring or LLM judging of the answer against the retrieved
documents~\citep{honovich2022true,laban2022summac,es2024ragas},
occasionally complemented by human annotation. Our frozen artifact uses
three automated support proxies: a DeBERTa-v$3$ zero-shot entailment
label score (\texttt{cross-encoder/nli-deberta-v3-base}), a
\texttt{roberta-large-mnli} zero-shot proxy, and a context-conditioned
RAGAS-style local judge. The supplement reports the exact scorer input
formats. Cross-scorer disagreement on the same fixed generations is
itself one of the strongest pieces of evidence for ControlledRAG
(\S\ref{sec:metric_fragility}).

\paragraph{RAG evaluation frameworks.} Recent systems such as
ARES~\citep{saadfalcon2024ares}, RAGChecker~\citep{ru2024ragchecker},
RAGTruth~\citep{niu2024ragtruth}, MIRAGE~\citep{park2025mirage},
FaithJudge~\citep{tamber2025faithjudge}, RAGAS~\citep{es2024ragas},
TRUE~\citep{honovich2022true}, and SummaC~\citep{laban2022summac}
provide evaluators, datasets, benchmark suites, or diagnostic metrics.
ControlledRAG is complementary. It does not replace those metrics or
leaderboards; it asks whether a RAG faithfulness claim survives the
hidden degrees of freedom those tools often instantiate.

\paragraph{Coherence and context structure.} Discourse coherence has a
long history in classical text
generation~\citep{li2014coherence,mesgar2018neural,nguyen2021conditional}
but, to our knowledge, has not been quantified as a deployment-time
RAG signal. \citet{shi2023large} show that LLMs are distracted by
irrelevant context but do not consider the case where each individual
passage is relevant yet the \emph{set} is internally fragmented;
\citet{liu2024lostmiddle} document position-dependent context
utilisation in long contexts. We treat the Context Coherence Score
(CCS) and answer-bearing evidence presence as two complementary
context-structure variables to \emph{report} alongside per-passage
similarity, not as substitutes for human evaluation or for cross-scorer
checks.

\paragraph{What we add.} Existing work proposes new retrievers,
generators, scorers, datasets, or auxiliary models. We do not.
ControlledRAG is not another metric suite or leaderboard. It is a
disclosure standard: it asks authors to show which generator, retriever,
context-structure, scorer, threshold, human-calibration, and cost
choices their claim survives.

\section{Setup and Audited Cells}
\label{sec:setup}

We hold the entire RAG pipeline fixed across audit experiments and vary
only the audit axis under test. The audited cells are deliberately
narrow---primarily short-answer extractive QA on SQuAD and PubMedQA
with $7$B-class generators---so that scorer, retriever, and generator
choices can be cleanly separated.

\paragraph{Datasets.} The matched-context, scaled, metric-fragility,
and span-control experiments use SQuAD~\citep{rajpurkar2016squad}; the
threshold-transfer and noise experiments additionally use
PubMedQA~\citep{jin2019pubmedqa}, HotpotQA~\citep{yang2018hotpotqa},
NaturalQuestions~\citep{kwiatkowski2019nq}, and TriviaQA. The long-form
stress test on QASPER and MS-MARCO is supplement-only and not treated
as broad generalization evidence (\S\ref{sec:limits} and supplement).

\paragraph{Pipeline.} Embeddings come from
\texttt{sentence-transformers/all-MiniLM-L$6$-v$2$}; indexing uses
ChromaDB; refinement uses
\texttt{cross-encoder/ms-marco-MiniLM-L-$6$-v$2$}. Default chunk size is
$1024$ characters with $10\%$ overlap and default top-$k$ is $3$.
Generators are served via local Ollama (Mistral-$7$B,
Qwen$2.5$-$7$B). The matched-context audit also re-encodes the same
$200$ pairs with BGE-large and E5-large.

\paragraph{Conditions.} Three retrieval configurations are evaluated
head-to-head per query: \textbf{baseline} (top-$k$ dense retrieval, no
refinement), \textbf{HCPC-v$1$} (Hybrid Context-Preserving Chunking:
refine all retrieved passages through sub-chunking and cross-encoder
selection---the original failure-mode condition), and
\textbf{HCPC-v$2$} (a coherence-gated HCPC variant with top-rank
protection, a refinement cap, and merge-back). The audit treats these
three conditions as instruments to measure the audit surface, not as
systems to crown a winner.

\paragraph{Metrics.} Per-query \textbf{faithfulness} in the legacy
audit tables means the stored automated support score in $[0,1]$;
\textbf{hallucination} is the indicator that this score is $<0.5$
unless a table explicitly states another label source. The frozen
DeBERTa score averages sentence-level
\texttt{cross-encoder/nli-deberta-v3-base} zero-shot entailment-label
scores over generated answer sentences. The second automated score is a
\texttt{roberta-large-mnli} zero-shot entailment-label proxy. In both
legacy NLI-proxy implementations, the answer sentence is the sequence
classified against the labels entailment/neutral/contradiction with
\texttt{hypothesis\_template="\{\}"}; the retrieved context remains in
the row but is not consumed by that proxy. The RAGAS-style judge is
context-conditioned: local \texttt{Mistral-$7$B} receives the question,
retrieved context, and answer and returns a support score on $[0,1]$
with a short rationale. This implementation detail is precisely why
scorer choice is an audit axis rather than a hidden detail.
\textbf{Retrieval similarity} is the mean cosine similarity between the
query and each retained chunk. \textbf{CCS} is the mean minus the
standard deviation of off-diagonal pairwise cosine similarities between
retained chunk embeddings under
\texttt{sentence-transformers/all-MiniLM-L6-v2}. The answer-span
indicator is whether any normalised gold-answer alias appears in the
retrieved context (lowercased, punctuation-stripped,
whitespace-normalised lexical match, no model).

\paragraph{Statistical protocol.} The matched-context test is
pre-registered with paired Wilcoxon and a $10{,}000$-resample
bootstrap CI on the mean paired difference; the scaled audit reports
seed-level paired Wilcoxon over five seeds; binary rates are reported
with Wilson $95\%$ CIs; the answer-span/control models report bootstrap
$95\%$ CIs on $R^2$ and AUROC. P-values outside the pre-registered
primary matched tests are descriptive across multiple audit cells; we
do not treat exploratory $p<0.05$ rows as standalone discoveries. All
audit cells, paired contrasts, and per-query CSVs are included in the
anonymized review artifact.

\section{ControlledRAG: The Minimum Reporting Standard}
\label{sec:controlledrag}

The seven-axis executive summary appeared in \S\ref{sec:intro}. This
section turns it into a reporting standard. For each axis, the author
must report what was varied or held fixed, why the axis matters, and
which overclaim the disclosure prevents.

\begin{table}[!htb]
\centering
\caption{ControlledRAG seven-axis audit map. The table is a reporting
standard, not a method prescription.}
\label{tab:controlledrag_audit_map}
\scriptsize
\setlength{\tabcolsep}{2pt}
\begin{tabular}{@{}p{0.12\linewidth}p{0.21\linewidth}p{0.21\linewidth}p{0.18\linewidth}p{0.17\linewidth}@{}}
\toprule
Axis & What to report & Why it matters & Evidence here & Overclaim prevented \\
\midrule
Generator & model family, size, decoding, fixed vs regenerated outputs &
same context can yield generator-specific behavior & Mistral main audit; Qwen$2.5$ compact replication &
generator-universal claims \\
Retriever & embedder, index, reranker, refinement policy &
retriever identity changes the context surface & MiniLM plus BGE/E5/GTE sanity checks &
``retrieval quality'' as one variable \\
Context structure & CCS, answer span, similarity, retained chunks &
support and coherence can diverge & matched CCS null; span/control diagnostic &
CCS as oracle or causal mediator \\
Scorer & scorer names, input format, thresholds, correlations &
same generations can produce scorer-dependent contrasts & three-scorer fixed-output audit &
single-score faithfulness claims \\
Human calibration & label protocol, raters, agreement, scorer alignment &
automated scores are not reference labels & $n=99$ two-rater calibration &
metric ranking as settled \\
Threshold & tuned $\tau$, tune data, off-diagonal transfer &
tuned policies may not transport & $5\times5$ transfer table &
reusable threshold claims \\
Cost & latency, indexing time, external-call assumptions, hardware &
deployment rankings depend on resource costs & CRAG / HCPC-v$2$ / RAPTOR-$2$L &
leaderboard without cost \\
\bottomrule
\end{tabular}
\end{table}

\begin{enumerate}
  \item \textbf{Generator.} Report the generator family, parameter
        count, decoding settings, and whether outputs are fixed before
        scoring or regenerated per condition. This prevents
        Mistral-only conclusions from being read as generator-general
        claims; the Qwen$2.5$ replication is a compact example, not a
        complete generator panel.
  \item \textbf{Retriever.} Report the embedding model, indexer,
        reranker, and any post-retrieval refinement step. This prevents
        MiniLM/BGE/E5 differences from being hidden under the word
        ``retrieval''; stronger embeddings can change the retrieval
        surface (\S\ref{sec:matched_ccs} and supplement).
  \item \textbf{Context structure.} Report at least two context-set
        diagnostics beyond per-passage similarity. CCS, mean query
        similarity, and answer-span presence are jointly informative
        (\S\ref{sec:matched_ccs}); this prevents coherence, similarity,
        and answer support from being collapsed into one retrieval
        quality number.
  \item \textbf{Scorer.} Report at least two faithfulness scorers, their
        input format, and a Pearson or Spearman cross-scorer
        correlation. Raw score contrasts and standardized contrasts are
        scorer-sensitive on the same fixed generations
        (\S\ref{sec:metric_fragility}); this prevents single-metric
        faithfulness claims.
  \item \textbf{Human calibration.} Report a multi-rater reference human
        label set, inter-rater agreement, and scorer correlation against
        the reference labels (\S\ref{sec:human_cal}). Calibration is not
        scorer endorsement; it prevents automatic metrics from being
        treated as reference labels.
  \item \textbf{Threshold transfer.} Report how any tuned threshold
        $\tau$ transfers across datasets. Cross-dataset transfer is
        uneven (\S\ref{sec:threshold_transfer}); this prevents a tuned
        threshold from being presented as a reusable policy without
        off-diagonal evidence.
  \item \textbf{Cost.} Report latency, indexing time, and external-call
        cost for every condition, including baselines. This prevents
        method rankings that ignore deployment cost
        (\S\ref{sec:cost}).
\end{enumerate}

The contribution is the standard, not any specific number. The rest of
the paper supplies one matched, pre-registered self-application in
which each empirical result motivates a required disclosure axis.

\section{The Matched-Context Audit}
\label{sec:matched_ccs}

The matched-context audit asks the diagnostic question directly:
\emph{at fixed mean per-passage query--passage similarity, does higher
CCS predict higher faithfulness?} The null answer is not the end of the
analysis. It creates a constructive diagnostic question: what variables
distinguish answer support once mean similarity is fixed?

\paragraph{Construction.} For SQuAD with Mistral, we form $200$
HIGH/LOW pairs from the per-query audit data so that mean
query--passage similarity is matched to within $0.02$ between members
of a pair while CCS is large in HIGH and small in LOW. Mean similarity
is $0.4236$/$0.4215$ across HIGH/LOW; mean CCS is $0.670$/$0.137$. The
similarity Wilcoxon is significant only because the residual noise is
two orders of magnitude smaller than the CCS gap, not because the
similarity gap is meaningful in absolute terms.

\paragraph{Result.} Table~\ref{tab:matched_ccs_main} reports the
pre-registered HIGH--LOW contrast. The mean paired difference is
$-0.002$, paired Wilcoxon $p = 0.628$, Cohen's $d_z = -0.017$,
$10{,}000$-resample bootstrap $95\%$ CI $[-0.022, +0.017]$. The HIGH
hallucination rate is also \emph{higher} than the LOW rate
($16.5\%$ vs.\ $9.0\%$). The pre-registered $H_1$
(HIGH$>$LOW faithfulness) is not supported in the audited cell. The
continuous mean contrast and the thresholded hallucination contrast
measure different parts of the score distribution, so a null mean can
coexist with threshold-sensitive tail differences. The matched test
rules out the simple HIGH-CCS-improves-faithfulness story in this
audited cell.

\begin{table}[!htb]
\centering
\caption{Matched-context audit on SQuAD/Mistral ($n{=}200$ pairs,
matched mean query--passage similarity).}
\label{tab:matched_ccs_main}
\footnotesize
\begin{tabular}{lcccc}
\toprule
Quantity & HIGH-CCS & LOW-CCS & Difference & Test \\
\midrule
Mean similarity & $0.4236$ & $0.4215$ & $+0.0021$ & matched by construction \\
Mean CCS        & $0.670$  & $0.137$  & $+0.533$  & --- \\
Faithfulness    & $0.636$  & $0.639$  & $-0.002$  & paired Wilcoxon $p=0.628$ \\
Hallucination   & $16.5\%$ & $9.0\%$  & $+7.5$ pp & McNemar $p=0.011$ \\
$95\%$ CI on faith diff & --- & --- & $[-0.022, +0.017]$ & $10{,}000$-resample bootstrap \\
\bottomrule
\end{tabular}
\end{table}

\paragraph{Span presence in matched contexts.} The first candidate is
answer-bearing evidence.
The lexical span detector finds the gold answer string more often in
LOW-CCS contexts than in HIGH-CCS contexts ($27.0\%$ vs.\ $17.0\%$,
McNemar $p=0.011$). This is itself a useful finding: under the
matched-similarity construction, the structural property captured by
CCS is not the same property as ``the retrieved context contains the
answer''. The lexical span variable is also a weak proxy: a context can
contain the answer string while failing to support the generated answer,
and thresholded hallucination and mean score can diverge. ControlledRAG
does not discard context structure; it decomposes it.

\paragraph{Answer-span/control diagnostic.}
Table~\ref{tab:span_control_main} extends the span check into a
regression and classification panel on the matched cell ($n=400$
contexts). We fit single-variable, two-variable, and three-variable
models against (i) continuous DeBERTa faithfulness ($R^2$) and (ii) the
binary hallucination indicator (AUROC).

\begin{table}[!htb]
\centering
\caption{Answer-span/control diagnostic on the matched SQuAD/Mistral
cell. $R^2$ is for continuous DeBERTa faithfulness; AUROC is for the
binary hallucination label. Bootstrap $95\%$ CIs are in the
supplement.}
\label{tab:span_control_main}
\footnotesize
\begin{tabular}{lcc}
\toprule
Predictors & Faithfulness $R^2$ & Hallucination AUROC \\
\midrule
answer span present                       & $0.019$ & $0.531$ \\
mean query similarity                     & $0.002$ & $0.586$ \\
CCS                                       & $0.003$ & $0.639$ \\
answer span $+$ similarity                & $0.024$ & $0.590$ \\
answer span $+$ CCS                       & $0.022$ & $0.631$ \\
answer span $+$ similarity $+$ CCS        & $\mathbf{0.025}$ & $\mathbf{0.643}$ \\
\bottomrule
\end{tabular}
\end{table}

\paragraph{Interpretation.} Three patterns are robust within the
audited cell. (i) For continuous faithfulness, answer-span presence is
the strongest \emph{single} predictor. (ii) For binary hallucination,
CCS is the strongest single predictor; the lexical span check is barely
above chance. (iii) The full feature set has the highest score on both quantities.
The $R^2$ values are small in absolute terms, so the table is not an
explanatory model of faithfulness. Its value is comparative: answer
support and context structure diagnose complementary failure modes, and
a faithfulness claim is more informative when both are reported. We
deliberately do not claim that either variable has a directional role, that
span subsumes CCS, or that CCS subsumes span.

\section{The Scaled Audit}
\label{sec:scaling}

The matched-context audit is intentionally local. The scaled audit
asks the same headline question---does HCPC-v$1$ refinement reliably
hurt faithfulness on SQuAD/Mistral?---at $n{=}2{,}500$ examples per
condition with five seeds, with each seed analysed by paired Wilcoxon
on $500$-pair contrasts.

\paragraph{Result.} Table~\ref{tab:scaling_main} reports the headline
table over the five seeds. The pooled means are
$0.661 / 0.650 / 0.661$ for baseline / HCPC-v$1$ / HCPC-v$2$, with
overlapping Wilson CIs on hallucination rate ($14.6\% / 14.7\% /
14.3\%$). Per-seed paired Wilcoxon yields one of five seeds significant
at $p<0.05$ for the baseline$-$HCPC-v$1$ contrast (seed 44, paired
$p=0.008$); the other four seeds yield $p \in [0.083, 0.234]$. The
v$2$$-$v$1$ contrast reaches significance in $3$ of $5$ seeds, but
the magnitude is small (per-seed mean diff $0.003$--$0.019$).

\begin{table}[!htb]
\centering
\caption{Scaled audit on SQuAD/Mistral, $n=2{,}500$ per condition,
pooled across five seeds.}
\label{tab:scaling_main}
\footnotesize
\begin{tabular}{lcccc}
\toprule
Condition & Faithfulness ($95\%$ CI) & Hallucination (Wilson) & Retrieval sim.\ & Refine rate \\
\midrule
Baseline   & $0.661 \, [0.655, 0.667]$ & $14.6\% \, [13.3, 16.0]$ & $0.532$ & $0.000$ \\
HCPC-v$1$  & $0.650 \, [0.645, 0.656]$ & $14.7\% \, [13.4, 16.2]$ & $0.569$ & $1.000$ \\
HCPC-v$2$  & $0.661 \, [0.655, 0.667]$ & $14.3\% \, [13.0, 15.7]$ & $0.532$ & $0.705$ \\
\bottomrule
\end{tabular}
\end{table}

\paragraph{Interpretation.} The original $n=30$ headline contrast
($0.7995 \to 0.6407$) does not survive at scale. The scaled contrast on
SQuAD/Mistral is small and seed-sensitive, and the headline should not
be cited as a broad refinement-hurts-faithfulness claim. This finding
is itself a ControlledRAG-style audit lesson: at $n=30$, the contrast
looked clean enough to be reported as a phenomenon, but the contrast at
$n=2{,}500$ rejects that reading without further controls. We treat the
remaining audits as the controls, not as reasons to reverse the
original claim.

\section{Fixed-Generation Metric Fragility}
\label{sec:metric_fragility}

This section is the core empirical pillar for the scorer axis. For the
scaled audit's $7{,}500$ Mistral rows, outputs are fixed before
rescoring, retrieval condition is fixed, and the only changed component
is the scorer. The three scorers are not guaranteed to be calibrated to
the same numerical scale, so raw score differences are not interpreted
as directly comparable causal estimates. The important claim is that
the sign, magnitude, and ranking of the reported contrast are
scorer-sensitive even on fixed generations. We then replicate the raw
contrast compactly on a second generator (Qwen$2.5$) at $n=100$ fixed
contexts per condition.

\begin{table}[!htb]
\centering
\caption{Fixed-generation metric fragility. Raw values are
baseline$-$HCPC-v$1$ score differences on each scorer's native scale.
Mistral z/rank columns standardize each scorer over the same fixed rows
before computing the paired contrast; bootstrap CIs are in the
supplement.}
\label{tab:metric_fragility_main}
\footnotesize
\begin{tabular}{lcccc}
\toprule
Scorer & Mistral raw & Mistral z & Mistral rank & Qwen$2.5$ raw \\
\midrule
DeBERTa proxy & $0.011$ & $0.071$ SD & $0.022$ & $0.015$ \\
roberta-large-MNLI proxy & $0.032$ & $0.231$ SD & $0.062$ & $0.046$ \\
RAGAS-style judge & $0.140$ & $0.337$ SD & $0.085$ & $0.134$ \\
\bottomrule
\end{tabular}
\end{table}

\begin{table}[!htb]
\centering
\caption{Cross-scorer correlations on fixed-generation rows.}
\label{tab:metric_correlations_main}
\footnotesize
\begin{tabular}{lcccc}
\toprule
Metric pair & Mistral $r$ & Mistral $\rho$ & Qwen$2.5$ $r$ & Qwen$2.5$ $\rho$ \\
\midrule
DeBERTa vs.\ roberta-large-MNLI       & $0.259$ & $0.265$ & $0.380$ & $0.429$ \\
DeBERTa vs.\ RAGAS-style              & $0.182$ & $0.212$ & $0.280$ & $0.318$ \\
roberta-large-MNLI vs.\ RAGAS-style   & $0.674$ & $0.651$ & $0.721$ & $0.668$ \\
\bottomrule
\end{tabular}
\end{table}

\paragraph{Scorer implementations.} The DeBERTa and
\texttt{roberta-large-mnli} rows are legacy zero-shot entailment-label
proxies from the frozen artifact; the local RAGAS-style judge is the
context-conditioned scorer. This is not a critique of any metric. It is
the audit point: scorer model, prompt, input formatting, and threshold
are part of the reported claim.

\paragraph{Interpretation.} Scorer choice is, on this audit, the single
largest movable quantity in the headline number. The human-calibration
panel in the next section also shows that no scorer should be treated as
an unqualified reference standard. Scorer choice is therefore an audit axis,
not an implementation detail: ControlledRAG requires at least two
scorers and a cross-scorer correlation to be reported. Scorer
triangulation is not proof of correctness; it discloses how much of the
reported number is scorer-dependent.

\section{Human Calibration}
\label{sec:human_cal}

Two raters (one expert, one general-domain) independently label $99$
matched-cell items with a faithful/hallucinated decision and a
$0$--$1$ confidence score. We use rater~$1$'s labels as the reference
human label set and rater~$2$'s labels as an independent reliability
check; we do not collapse the two raters by majority or discussion
(with two raters either procedure would be ill-defined). The supplement
details the protocol.

\paragraph{Agreement and correlation.} Raters reach Cohen's
$\kappa = 0.774$ on the binary label, which is in the substantial-
to high-agreement range for this kind of task. Spearman correlations
between the rater-$1$ label set and the three automated scorers are
$\rho_{\textrm{DeBERTa}} = 0.140$,
$\rho_{\textrm{second NLI}} = 0.380$, and
$\rho_{\textrm{RAGAS-style}} = 0.441$. Bootstrap $95\%$ CIs on the
three correlations overlap, so the metric ranking is suggestive
rather than settled.

\paragraph{Interpretation.} Two human raters who reach substantial
agreement still give an uneven picture of scorer alignment. The right
reading is calibration, not scorer endorsement: the panel reports how
well each scorer aligns with the reference human label set on this
audited cell, and that alignment is uneven enough to merit explicit
disclosure. ControlledRAG asks for this calibration as a routine
reporting axis, not as a means of crowning one scorer.

\section{Threshold Transfer}
\label{sec:threshold_transfer}

A faithfulness claim that depends on a tuned operating threshold $\tau$
should report whether $\tau$ transfers across datasets. We sweep
$\tau \in \{0.3, 0.4, 0.5, 0.6, 0.7\}$ for a CCS-gate retrieval
intervention. For each tune dataset, $\tau^\ast$ is the threshold that
maximises recovery on that dataset alone; diagonal recovery evaluates
that $\tau^\ast$ on the same dataset, while off-diagonal recovery
evaluates it on the other datasets. Recovery is a dimensionless
normalised contrast,
\[
R_\tau(d)=
\frac{\mathrm{faith}_{\mathrm{gate},\tau}(d) -
      \mathrm{faith}_{\mathrm{HCPC}\text{-}v1}(d)}
     {\mathrm{faith}_{\mathrm{baseline}}(d) -
      \mathrm{faith}_{\mathrm{HCPC}\text{-}v1}(d)}.
\]
Values can exceed $1$ when the gate beats the baseline relative to
HCPC-v$1$, and they can be negative or sign-inverted when the
baseline$-$HCPC-v$1$ denominator is small or negative. We therefore use
the table as a transport diagnostic, not as a universal utility score.
A tuned threshold is not a reusable policy unless this off-diagonal
transfer is reported.

\begin{table}[!htb]
\centering
\caption{Threshold transfer summary. For each tune dataset, $\tau^\ast$
is the value of $\tau$ that maximises recovery on that dataset alone;
``diagonal recovery'' is recovery at $\tau^\ast$ on the same dataset;
``off-diagonal mean'' is the mean recovery at $\tau^\ast$ over the four
\emph{other} eval datasets. A large positive gap means the tuned
$\tau^\ast$ does not transport. The full per-cell transfer matrix is
in supplement \S 9.}
\label{tab:threshold_transfer_main}
\footnotesize
\begin{tabular}{lcrrr}
\toprule
$\tau$ tuned on & $\tau^\ast$ & Diagonal recovery & Off-diagonal mean & Gap (diag $-$ off-diag) \\
\midrule
PubMedQA  & $0.4$ & $+1.452$ & $+0.451$ & $+1.002$ \\
NaturalQs & $0.5$ & $+1.401$ & $+0.211$ & $+1.191$ \\
SQuAD     & $0.3$ & $+0.823$ & $-0.141$ & $+0.963$ \\
TriviaQA  & $0.4$ & $+0.464$ & $+0.698$ & $-0.234$ \\
HotpotQA  & $0.3$ & $-0.112$ & $+0.093$ & $-0.205$ \\
\bottomrule
\end{tabular}
\end{table}

\paragraph{Interpretation.} PubMedQA, NaturalQuestions, and SQuAD show
a positive diagonal-vs-off-diagonal gap large enough to flag: tuning
$\tau$ on one of them and transporting it loses most of the
on-diagonal normalized contrast. TriviaQA and HotpotQA show a negative
gap; their $\tau^\ast$ does not win on their own data (HotpotQA's
diagonal value is itself negative), so transport is moot. The
supplement reports the full $5 \times 5$ per-cell matrix. ControlledRAG
asks for the transfer table to be reported because per-cell normalized
contrasts can change sign depending on which dataset $\tau$ is tuned
on; we do not claim universal threshold failure from this one sweep.

\section{Coherence vs.\ Generic Retrieval Noise}
\label{sec:noise}

A reasonable objection to any context-structure diagnostic is whether it
is just a re-statement of generic retrieval noise. We test this by
injecting controlled retrieval noise at three rates ($1/3$, $2/3$, and
fully replaced) and comparing two noise types: random off-topic
passages, and coherence-preserving uninformative passages (passages
drawn from the same topic distribution that do not contain the answer).

\paragraph{Result.} The faithfulness slope per noise rate is
$-0.069$ for random noise and $-0.043$ for coherence-preserving
uninformative noise. The similarity slope is much larger for random
noise ($-0.481$ per noise rate) than for coherence-preserving noise
($-0.113$), as expected. Coherence-preserving noise has a
distinguishably smaller faithfulness drop at the same nominal
substitution rate.

\paragraph{Interpretation.} On this matched comparison, coherence
carries some signal that is not captured by mean per-passage
similarity. We do not claim that the gap is large in absolute terms,
only that the two noise types produce distinguishable faithfulness
slopes at matched substitution rate. The right reading is that
coherence and similarity are correlated but not identical; ControlledRAG
asks for both to be reported when a context-structure claim is made.

\section{Cost-Aware Baselines}
\label{sec:cost}

Faithfulness numbers reported without latency and indexing cost are
under-identified along the cost axis. This section is not a leaderboard;
it is an identifiability test for deployment claims. We re-run a
head-to-head on $n=200$ SQuAD and $n=200$ HotpotQA queries with three
baselines: CRAG~\citep{yan2024crag}, HCPC-v$2$, and a $2$-level
RAPTOR~\citep{raptor2024} index.

\begin{table}[!htb]
\centering
\caption{Cost-aware head-to-head, $n=200$ per (dataset, condition).
Latency is end-to-end per query. Indexing time is the offline cost of
building the dataset's retrieval structure. Faithfulness CIs are
bootstrap $95\%$ intervals; hallucination CIs are Wilson $95\%$
intervals. RAPTOR-$2$L's indexing cost is $17$--$22\times$ the dense
baseline's.}
\label{tab:cost_main}
\scriptsize
\begin{tabular}{llcccc}
\toprule
Dataset & Condition & Faithfulness & Hallucination & Mean latency & Indexing \\
\midrule
SQuAD    & CRAG        & $0.698\,[0.675,0.722]$ & $12.0\%\,[8.2,17.2]$ & $1{,}440$ ms & $\phantom{1}8.0$ s \\
SQuAD    & HCPC-v$2$   & $0.708\,[0.685,0.732]$ & $12.5\%\,[8.6,17.8]$ & $1{,}720$ ms & $\phantom{1}8.0$ s \\
SQuAD    & RAPTOR-$2$L & $0.710\,[0.685,0.735]$ & $12.5\%\,[8.6,17.8]$ & $1{,}714$ ms & $175.4$ s \\
HotpotQA & CRAG        & $0.643\,[0.627,0.658]$ & $10.5\%\,[7.0,15.5]$ & $1{,}941$ ms & $\phantom{1}8.0$ s \\
HotpotQA & HCPC-v$2$   & $0.633\,[0.617,0.650]$ & $12.5\%\,[8.6,17.8]$ & $2{,}289$ ms & $\phantom{1}8.0$ s \\
HotpotQA & RAPTOR-$2$L & $0.631\,[0.615,0.647]$ & $13.0\%\,[9.0,18.4]$ & $2{,}414$ ms & $134.0$ s \\
\bottomrule
\end{tabular}
\end{table}

\paragraph{Harness and cost accounting.} This is a closed, zero-dollar
local/free-tier harness: local Ollama/Mistral generation, Chroma
indexing, MiniLM cross-encoders, and no paid APIs. The CRAG
reimplementation uses the documented evaluator/strip-refinement
structure with live web search disabled, so its external-call count and
external dollar cost are both zero in this cell. A live-web CRAG run
would need to report web-call count, API cost, and network latency as a
separate deployment setting.

\paragraph{Self-RAG.} Self-RAG~\citep{asai2024selfrag} is reported in
the supplement because its reflective decoding protocol violates the
fixed-context assumption used in the main audit. ControlledRAG asks
that any baseline of this kind be reported with its harness and
excluded from head-to-head averages it cannot fairly support.

\paragraph{Interpretation.} On SQuAD, all three systems are within
$1.2$ percentage points on faithfulness with overlapping intervals;
HCPC-v$2$ and RAPTOR-$2$L have the same hallucination rate while CRAG
is numerically lower. On HotpotQA, CRAG is numerically strongest on
faithfulness, hallucination, and latency in this small cell. RAPTOR-$2$L's offline
indexing time is $17$--$22\times$ the dense baseline's. There is no
method that is uniformly preferred once faithfulness, hallucination, latency, and
indexing cost are considered, and ControlledRAG's cost axis exists
precisely so that contrasts of this kind are reported rather than
chosen.

\section{Discussion}
\label{sec:discussion}

\paragraph{What survives.} The under-identification thesis is not a
collection of negative results; it is the reason a reporting standard is
needed. The matched intervention is null, the scaled headline is small
and seed-sensitive, fixed-generation metric fragility changes raw and
standardized reported contrasts, the threshold table has sign-changing
cells, and
the cost-aware head-to-head has no dominator. Each cell corresponds to a
ControlledRAG axis. The positive contribution is a standard of evidence:
RAG faithfulness claims should not be accepted from a single retrieval
score, faithfulness scorer, tuned threshold, or favorable baseline row.

\paragraph{What CCS is and is not.} CCS is one of several context-set
diagnostics ControlledRAG asks for. It is not a faithfulness oracle,
not a guarantee, and not a substitute for human evaluation or for
cross-scorer triangulation. The matched-context audit rules out a
clean directional mediating role for CCS in the audited cell. The
answer-span/control diagnostic shows that span and CCS are
complementary axes, with span carrying more continuous-faithfulness
signal and CCS carrying more hallucination-AUROC signal. We report
both; ControlledRAG asks readers to do likewise.

\paragraph{Where the audit stops short.} Two visible gaps remain. The
second-generator replication uses one local Qwen$2.5$ model, not a
broad generator panel; the long-form stress test on QASPER and MS-MARCO
has only $40$ distinct questions and is supplement-only. We report both
gaps as part of the audit because the standard is meant to make scope
visible, not to hide it.

\section{Related Work, Continued}
\label{sec:related}

Beyond the RAG-evaluation context in \S\ref{sec:bg}, this paper sits
next to a small but growing literature on RAG \emph{evaluation
methodology}. \citet{honovich2022true} re-evaluate factual-consistency
evaluation broadly and conclude that no single NLI-based scorer is
universally preferred; ControlledRAG's metric-fragility axis is the
operational consequence of that conclusion for RAG.
\citet{es2024ragas} propose RAGAS as a metric suite; we use a
RAGAS-style judge as one of three scorers and treat its disagreement
with NLI scorers as evidence for the standard, not as a critique of
RAGAS.

\begin{table}[!htb]
\centering
\caption{Positioning against recent RAG evaluation work.}
\label{tab:related_positioning}
\scriptsize
\begin{tabular}{p{0.17\linewidth}p{0.36\linewidth}p{0.36\linewidth}}
\toprule
Work & Primary contribution & ControlledRAG relation \\
\midrule
ARES~\citep{saadfalcon2024ares} & trained judges for context relevance, answer faithfulness, and answer relevance with PPI calibration & complementary evaluator; ControlledRAG asks whether claims survive scorer/retriever/generator choices \\
RAGChecker~\citep{ru2024ragchecker} & fine-grained retrieval and generation diagnostics for RAG systems & complementary metric framework; ControlledRAG is the disclosure layer around such diagnostics \\
RAGTruth~\citep{niu2024ragtruth} & manually annotated hallucination corpus with word-level labels & complementary benchmark; ControlledRAG asks how benchmark conclusions transport across audit axes \\
MIRAGE~\citep{park2025mirage} & RAG benchmark with component-specific metrics and adaptability probes & complementary benchmark; ControlledRAG emphasizes reporting hidden degrees of freedom \\
FaithJudge~\citep{tamber2025faithjudge} & LLM-as-judge faithfulness benchmark and evolving leaderboard & complementary scorer/leaderboard; ControlledRAG requires scorer calibration and cross-scorer disclosure \\
RAGAS / TRUE / SummaC~\citep{es2024ragas,honovich2022true,laban2022summac} & reference-free or NLI-style factual-consistency metrics & metrics used as audit inputs, not treated as reference standards \\
\bottomrule
\end{tabular}
\end{table}

\section{Limitations and Ethics}
\label{sec:limits}

\paragraph{Audited surface.} The audited cells are mostly short-answer
extractive QA on SQuAD with $7$B-class generators. The matched-context
audit is on SQuAD only. The conclusions are scoped to that surface.
Stronger embeddings (BGE-large, E5-large) preserve the matched
HIGH$>$LOW CCS ordering but the fixed-generation result remains null
because the generations are unchanged; the supplement reports the
multi-retriever generation table as a sanity check, not as exhaustive
robustness.

\paragraph{Generators.} The second-generator replication uses one
local Qwen$2.5$ model at $n=100$ fixed contexts per condition. The
metric-fragility ordering is preserved, but ControlledRAG explicitly
asks for a broader generator panel before any universal-generator claim
can be made. We do not make that claim here.

\paragraph{Long-form scope.} The QASPER/MS-MARCO long-form stress test
has $40$ distinct questions and is included only as an exploratory
scope probe in the supplement. The short-answer paradox does not
generalise cleanly to that small long-form sample (MS-MARCO span-faith
paradox drop $-0.033$, QASPER $0.002$). This is reported as a scope
limit, not as broad long-form generalization evidence.

\paragraph{Human evaluation.} Two raters at $n=99$ is calibration, not
a settled scorer ranking. Bootstrap CIs on the Spearman correlations
overlap; the human ranking is suggestive only.

\paragraph{Self-RAG harness.} The Self-RAG~\citep{asai2024selfrag}
appendix baseline runs in a harness mismatched with the audit's
fixed-context protocol. We report the row in the supplement and
exclude it from head-to-head averages. We explicitly do not generalize
from this mismatched harness to Self-RAG overall.

\paragraph{Ethics.} The audit uses public datasets and free-tier or
local compute. No human-subjects data is collected beyond the two-rater
calibration panel, which is not personally identifiable and does not
involve sensitive content. The review artifact is anonymized; public
release links will follow acceptance.

\paragraph{Scope as a feature.} These limitations are precisely the
axes ControlledRAG asks future papers to disclose. Our evidence is
deliberately scoped, and the standard makes that scope explicit.

\section{Conclusion}
\label{sec:conclusion}

A RAG faithfulness claim should not be a single number tied to a single
scorer, retriever, generator, threshold, or baseline. The audit
experiments in this paper, taken together, show that under any one of
those collapsed axes the reported headline can move materially.
The positive contribution is a prescriptive minimum reporting standard,
ControlledRAG, that names the seven axes a faithfulness claim has to
report together, and demonstrates the standard by self-applying it to
our own prior contribution. ControlledRAG turns RAG faithfulness
evaluation from a single-score claim into a seven-axis auditable
evidence standard.
